% Options for packages loaded elsewhere
\PassOptionsToPackage{unicode}{hyperref}
\PassOptionsToPackage{hyphens}{url}
\documentclass[
  ignorenonframetext,
]{beamer}
\newif\ifbibliography
\usepackage{pgfpages}
\setbeamertemplate{caption}[numbered]
\setbeamertemplate{caption label separator}{: }
\setbeamercolor{caption name}{fg=normal text.fg}
\beamertemplatenavigationsymbolsempty
% remove section numbering
\setbeamertemplate{part page}{
  \centering
  \begin{beamercolorbox}[sep=16pt,center]{part title}
    \usebeamerfont{part title}\insertpart\par
  \end{beamercolorbox}
}
\setbeamertemplate{section page}{
  \centering
  \begin{beamercolorbox}[sep=12pt,center]{section title}
    \usebeamerfont{section title}\insertsection\par
  \end{beamercolorbox}
}
\setbeamertemplate{subsection page}{
  \centering
  \begin{beamercolorbox}[sep=8pt,center]{subsection title}
    \usebeamerfont{subsection title}\insertsubsection\par
  \end{beamercolorbox}
}
% Prevent slide breaks in the middle of a paragraph
\widowpenalties 1 10000
\raggedbottom
\AtBeginPart{
  \frame{\partpage}
}
\AtBeginSection{
  \ifbibliography
  \else
    \frame{\sectionpage}
  \fi
}
\AtBeginSubsection{
  \frame{\subsectionpage}
}
\usepackage{iftex}
\ifPDFTeX
  \usepackage[T1]{fontenc}
  \usepackage[utf8]{inputenc}
  \usepackage{textcomp} % provide euro and other symbols
\else % if luatex or xetex
  \usepackage{unicode-math} % this also loads fontspec
  \defaultfontfeatures{Scale=MatchLowercase}
  \defaultfontfeatures[\rmfamily]{Ligatures=TeX,Scale=1}
\fi
\usepackage{lmodern}
\ifPDFTeX\else
  % xetex/luatex font selection
\fi
% Use upquote if available, for straight quotes in verbatim environments
\IfFileExists{upquote.sty}{\usepackage{upquote}}{}
\IfFileExists{microtype.sty}{% use microtype if available
  \usepackage[]{microtype}
  \UseMicrotypeSet[protrusion]{basicmath} % disable protrusion for tt fonts
}{}
\makeatletter
\@ifundefined{KOMAClassName}{% if non-KOMA class
  \IfFileExists{parskip.sty}{%
    \usepackage{parskip}
  }{% else
    \setlength{\parindent}{0pt}
    \setlength{\parskip}{6pt plus 2pt minus 1pt}}
}{% if KOMA class
  \KOMAoptions{parskip=half}}
\makeatother
\usepackage{longtable,booktabs,array}
\usepackage{caption}
\captionsetup[table]{skip=6pt}
\newcounter{none} % for unnumbered tables
\usepackage{calc} % for calculating minipage widths
\usepackage{caption}
% Make caption package work with longtable
\makeatletter
\def\fnum@table{\tablename~\thetable}
\makeatother
\setlength{\emergencystretch}{3em} % prevent overfull lines
\providecommand{\tightlist}{%
  \setlength{\itemsep}{0pt}\setlength{\parskip}{0pt}}
% Auto-generated by infrastructure/rendering/_pdf_latex_helpers.py::extract_math_font_preamble.
% Minimal header passed to Pandoc via -H so Beamer slides pick up the same math font
% as the combined-PDF path. Adjust the manuscript's preamble.md to change the font globally.
\usepackage{unicode-math}
\setmathfont{latinmodern-math.otf}

% Manuscript macro/package fallbacks (newcommand -> providecommand).
% Core mathematics
\usepackage{amsmath}
\usepackage{amssymb}
% Document layout
% Tables
\usepackage{booktabs}
% Algorithm typesetting (for pseudocode in §2 Methodology)
% Code listings
\usepackage{listings}
% Typography and formatting
% Cross-references and citations
% ── Unicode-capable mono font for code listings ──────────────────────
% JuliaMono (TeX Live 2026) covers the full math/Greek glyph set used in
% scientific code blocks; the default lmmono lacks \alpha/\mu/\partial/\nabla.
%
% The hardcoded TeX Live 2026 path below is portable to most macOS / Linux
% TeX Live installs that placed JuliaMono under the default texmf-dist path.
% If your TeX Live version differs (e.g., 2025, 2027) or your distribution
% installs fonts elsewhere, comment out the explicit Path = line — fontspec
% will then search the system font path. If JuliaMono is not installed at
% all, comment out the whole block; xelatex falls back to lmmono with
% reduced Unicode coverage but the document still compiles.
% Math font for unicode-math: Latin Modern Math (TeX Live) has full BMP
% coverage including U+2223 (\mid), U+226A/226B (\ll/\gg), and the Greek/
% blackboard letters used in equations. Without an explicit \setmathfont,
% unicode-math falls back to lmroman text font which lacks several glyphs
% and emits "Missing character" warnings on every \mid in math mode.

% Slide layout defaults for warning-clean scientific decks.
\usepackage{etoolbox}
\IfFileExists{xurl.sty}{\usepackage{xurl}}{}
\IfFileExists{seqsplit.sty}{\usepackage{seqsplit}}{\newcommand{\seqsplit}[1]{#1}}
\protected\def\breakseq#1{\seqsplit{#1}}
\protected\def\breaktt#1{\begingroup\ttfamily\seqsplit{#1}\endgroup}
\setlength{\emergencystretch}{6em}
\tolerance=5000
\hbadness=10000
\hfuzz=1pt
\setlength{\tabcolsep}{2pt}
\AtBeginEnvironment{longtable}{\tiny\renewcommand{\arraystretch}{0.86}\setlength{\tabcolsep}{1pt}}
\AtBeginEnvironment{tabular}{\tiny\renewcommand{\arraystretch}{0.86}\setlength{\tabcolsep}{1pt}}
\AtBeginEnvironment{equation}{\tiny}
\AtBeginEnvironment{equation*}{\tiny}
\AtBeginEnvironment{align}{\tiny}
\AtBeginEnvironment{align*}{\tiny}
\AtBeginEnvironment{itemize}{\footnotesize}
\AtBeginEnvironment{enumerate}{\footnotesize}
\AtBeginEnvironment{description}{\footnotesize}
\setbeamerfont{caption}{size=\tiny}
\setbeamerfont{caption name}{size=\tiny}
\setbeamerfont{normal text}{size=\small}
\setbeamerfont{frametitle}{size=\small}
\setbeamerfont{section title}{size=\footnotesize}
\setbeamerfont{subsection title}{size=\footnotesize}
\setbeamertemplate{section page}{%
  \centering
  \begin{beamercolorbox}[sep=12pt,center,wd=\paperwidth]{section title}
    \parbox{0.86\paperwidth}{\centering\usebeamerfont{section title}\insertsection\par}
  \end{beamercolorbox}
}
\setbeamertemplate{subsection page}{%
  \centering
  \begin{beamercolorbox}[sep=8pt,center,wd=\paperwidth]{subsection title}
    \parbox{0.86\paperwidth}{\centering\usebeamerfont{subsection title}\insertsubsection\par}
  \end{beamercolorbox}
}
\setlength{\abovecaptionskip}{2pt}
\setlength{\belowcaptionskip}{0pt}

% Natbib and cross-reference fallbacks — slides don't load natbib
% or cleveref, but combined-PDF manuscript prose may emit these
% commands. The fallback renders citations as a bracketed key list
% and cross-references as detokenized labels so slides don't fail on
% undefined control sequences or raw underscores. \providecommand is
% a no-op if packages load later.
\providecommand{\citep}[1]{[#1]}
\providecommand{\citet}[1]{#1}
\providecommand{\citealp}[1]{#1}
\providecommand{\citeauthor}[1]{#1}
\providecommand{\citeyear}[1]{#1}
\providecommand{\cref}[1]{\texttt{\detokenize{#1}}}
\providecommand{\Cref}[1]{\texttt{\detokenize{#1}}}
\providecommand{\crefrange}[2]{\texttt{\detokenize{#1}}--\texttt{\detokenize{#2}}}
\providecommand{\Crefrange}[2]{\texttt{\detokenize{#1}}--\texttt{\detokenize{#2}}}
\providecommand{\paragraph}[1]{\textbf{#1}\ }

\newtheorem{proposition}{Proposition}
\newtheorem{hypothesis}{Hypothesis}
\newtheorem{remark}[theorem]{Remark}
\newtheorem{axiom}{Axiom}
\newtheorem{property}{Property}
\usepackage{bookmark}
\IfFileExists{xurl.sty}{\usepackage{xurl}}{} % add URL line breaks if available
\urlstyle{same}
% fallback for those not using the hyperref driver hyperxmp:
\makeatletter
\@ifundefined{xmpquote}{\newcommand{\xmpquote}[1]{#1}}{}
\makeatother
\hypersetup{
  hidelinks,
  pdfcreator={LaTeX via pandoc}}

\author{\texorpdfstring{}{}}
\date{}

\begin{document}

\section{Configuration Generation Is Not
Authorization}\label{sec:configuration_authorization}

\begin{frame}[allowframebreaks]{Configuration Generation Is Not
Authorization}
An AI agent can write a declarative configuration, an access-control
policy, or a deployment change as fluently as it writes code. Generation
is not the security event. The security event is whatever independently
enforces the boundary between what the agent may change and what it may
only request --- above all, the mechanisms that constrain the agent
itself. This section deepens that argument, states the 9 review
invariants that an agent must never be able to violate, and locates the
separation of policy writing, policy approval, and permission exercise
inside the authority architecture of sec.~10.
\end{frame}

\begin{frame}[allowframebreaks]{Generation versus enforcement}
\protect\phantomsection\label{generation-versus-enforcement}
The distinction is complete mediation (Saltzer and Schroeder 1975):
every access to a protected resource is checked by an enforcement point,
and the authorship of the request is irrelevant to whether the check
runs. A proposed configuration is a request. Declarative formats make
requests legible --- diffable, reviewable, testable --- which is why
they are attractive in this threat model; but legibility is a property
of the document, while enforcement is a property of the system that acts
on it. A policy nobody enforces is a wish. An enforced policy whose
enforcement mechanisms the agent can reconfigure is the same wish, one
step removed, and it is the more dangerous form because it advertises
control that does not exist.

\end{frame}

\begin{frame}[allowframebreaks]{Generation versus enforcement}
The authorized-misuse path of sec.~3 runs straight
through configuration. An agent need not disable isolation through an
exploit if it can propose a configuration generation that disables
isolation and have that generation applied as routine system management.
The attack is a policy review that did not happen, not a boundary that
was crossed.
\end{frame}

\begin{frame}[fragile,allowframebreaks]{Declarative configuration as an
inspectable policy surface}
\protect\phantomsection\label{declarative-configuration-as-an-inspectable-policy-surface}
Used correctly, declarative configuration is the strongest inspectable
policy surface in common use. The NixOS model makes intended system
state explicit, keeps generations diffable, and supports activation and
rollback of configuration as a first-class operation ({community} 2026a,
2026b). Review can compare what will be deployed against what was
approved; drift becomes visible; policy becomes data. These are real
gains over imperative mutation, and this review's recommendations in
sec.~6 lean on them.

\end{frame}

\begin{frame}[fragile,allowframebreaks]{Declarative configuration as an
inspectable policy surface}
Two traps sit alongside the gains. The first is inspection without
enforcement: reviewing the file changes nothing if the deployment
pipeline applies whatever appears next. The second is enforcement
without independence: a reproducibly deployed privilege escalation
remains a privilege escalation, and reproducibility makes it
deterministic and fleet-wide --- every machine receives the escalation
in the same rebuild. Reproducible deployment guarantees the same result
everywhere; it does not ask whether the result is authorized (Project
2026a). Credential handling shows the same trap from the secrets side:
declarative cleanliness that embeds secrets in configuration exposes
them to every store reader and possibly to external caches, so the
documented practice reads secrets at runtime under access control
(Project 2026b). Inspection value and authorization value are different
properties, and the second never follows from the first.
\end{frame}

\begin{frame}[fragile,allowframebreaks]{Declarative configuration as an
inspectable policy surface}

Provenance inspection, in particular, is maturing from a manual
discipline into packaged tooling, and the Nix ecosystem is where it is
furthest along. Source provenance is now a required attribute in
nixpkgs: a \breaktt{meta.sourceProvenance} declaration tagging each
package's inputs --- from official releases to unvendored source trees
--- is merged as a requirement rather than a convention, which makes the
provenance question answerable from the derivation data itself rather
than from documentation ({community} 2026c). Around that requirement a
tooling stack has formed: sbomnix generates CycloneDX and SPDX software
bills of materials for a NixOS closure and emits SLSA provenance
attestations for builds ((TII) 2026); Trustix rebuilds packages
independently and compares the results bit-for-bit against what the
binary cache distributed, converting trust in a substitute into a
checkable claim ({community} 2026d); and the
\end{frame}

\begin{frame}[fragile,allowframebreaks]{Declarative configuration as an
inspectable policy surface}
\breaktt{nix\ provenance\ show} and \breaktt{nix\ provenance\ verify}
subcommands packaged by Determinate Systems expose the provenance of a
built artifact for inspection and verification as ordinary command-line
operations. The direction is exactly right for this threat model: the
provenance of deployed software becomes inspectable data, in the same
way declarative configuration made intended state inspectable. The
section's caution nonetheless applies unchanged. Provenance attests how
an artifact was built and from what inputs; it does not authorize
deployment, and a perfectly attested artifact is not thereby a permitted
one. Inspection tooling raises the quality of the review; the
authorization decision remains with the enforcement points this section
describes.
\end{frame}

\begin{frame}[fragile,allowframebreaks]{The review invariants}
\protect\phantomsection\label{the-review-invariants}
tbl.~\ref{tbl:invariants} states the review invariants: what an agent
must never be able to do, the failure each invariant closes, and where
enforcement must live.

\end{frame}

\begin{frame}[fragile,allowframebreaks]{The review invariants}
\begin{longtable}[]{@{}
  >{\raggedright\arraybackslash}p{(\linewidth - 4\tabcolsep) * \real{0.3333}}
  >{\raggedright\arraybackslash}p{(\linewidth - 4\tabcolsep) * \real{0.3333}}
  >{\raggedright\arraybackslash}p{(\linewidth - 4\tabcolsep) * \real{0.3333}}@{}}
\caption{The 9 configuration review invariants: the capability closed,
the failure each prevents, and the enforcement
point}\label{tbl:invariants}\tabularnewline
\toprule\noalign{}
\begin{minipage}[b]{\linewidth}\raggedright
Invariant: the agent must never be able to
\end{minipage} & \begin{minipage}[b]{\linewidth}\raggedright
Failure the invariant closes
\end{minipage} & \begin{minipage}[b]{\linewidth}\raggedright
Enforcement point
\end{minipage} \\
\midrule\noalign{}
\endfirsthead
\toprule\noalign{}
\begin{minipage}[b]{\linewidth}\raggedright
Invariant: the agent must never be able to
\end{minipage} & \begin{minipage}[b]{\linewidth}\raggedright
Failure the invariant closes
\end{minipage} & \begin{minipage}[b]{\linewidth}\raggedright
Enforcement point
\end{minipage} \\
\midrule\noalign{}
\endhead
Add an unrestricted credential provider & Every task gaining reach into
secrets through one configuration edit & Credential mediation service;
provider registry outside agent-writable configuration \\
Disable isolation & A prevention failure becoming whole-machine
compromise by policy rather than exploit & Isolation baseline outside
agent-writable generations; changes require external approval and a
verified boot path \\
Widen its own network policy & The egress boundary rewritten toward
attacker-chosen destinations, reopening data-loss paths
(sec.~12) & Egress policy owned by a boundary the agent
cannot rewrite \\
Replace the trusted update signer & Supply-chain takeover achieved as
ordinary configuration change & Signing keys and trust anchors held
outside every agent-writable generation \\
Alter the approving identity & An attacker redefining who counts as an
approver, then self-approving & Approving-identity registry immutable
from any agent context \\
Mount unrelated user data & Cross-domain reads into
\breaktt{personal\_identity} or \breaktt{credential\_service} data & Mount
policy scoped per trust domain (tbl.~8) \\
Expand its own tool grants & Self-issued authority through the
tool-bridge surface (sec.~12) & Grant registry held and
reviewed as explicit authority grants \\
Rewrite the audit trail & Destruction of the record incident response
depends on & Append-only audit storage outside every execution
environment (tbl.~9) \\
Approve its own policy changes & Closure of the loop from proposal to
authorization to exercise & Write-policy / approve-policy / exercise
separation (below) \\
\bottomrule\noalign{}
\end{longtable}

\end{frame}

\begin{frame}[fragile,allowframebreaks]{The review invariants}
The 9 invariants share one design principle: the state that would relax
an invariant is never reachable from agent-writable configuration. They
are capability invariants, not instruction invariants. A well-prompted
agent and a compromised agent must be equally unable to violate them,
because the authorized-misuse path does not distinguish persuasion from
compromise at the enforcement layer. An invariant that holds only while
the agent behaves is not an invariant; it is a preference.
\end{frame}

\begin{frame}[allowframebreaks]{Separating policy writing, policy
approval, and permission exercise}
\protect\phantomsection\label{separating-policy-writing-policy-approval-and-permission-exercise}
The write-policy / approve-policy / exercise-permission separation maps
onto the authority ladder of sec.~10: proposal
and staging are distinct rungs from authorization, and exercise is
distinct from both. The operating rule is that no single principal holds
more than one of the three roles for a consequential change. The
identity that writes a policy change differs from the identity that
approves it, and whatever exercises the new permission does so only
after an approval that neither writer nor exerciser controls. Hardy's
confused deputy is precisely the closure of this loop --- the program
that compiles the request and holds the authority over the resource it
names (Hardy 1988) --- and the separation is the structural answer:
split the deputy into roles that cannot be merged by the party under
\end{frame}

\begin{frame}[allowframebreaks]{Separating policy writing, policy
approval, and permission exercise}
attack.

Deterministic policy checks can automate approval where the check is
complete: binary, fully specified, and independent of the task's
success. Where judgment is required, the approver receives the full
operation context --- artifact, destination, scope, expiry --- not a
summary, mirroring the approval-context discipline of
sec.~13. An approver who sees ``network policy
update'' has approved nothing; an approver who sees the exact diff that
widens egress to a named destination has at least been given the choice.

\end{frame}

\begin{frame}[allowframebreaks]{Separating policy writing, policy
approval, and permission exercise}
Signer management is the sharpest instance of the separation. A
configuration change that replaces the trusted update signer --- or the
signer of a binary cache --- converts every subsequent update into
attacker-chosen code, so the review must ask who holds signing authority
rather than merely whether signatures verify; trusting a signer is not
proof that an artifact was produced by the expected derivation
({community} 2026e). Invariant 4 in tbl.~\ref{tbl:invariants} is
therefore enforced by construction: signing trust anchors live outside
every agent-writable generation, and changing them is the
highest-consequence generation request an agent can make --- which is
exactly why it must be one the agent cannot complete alone.
\end{frame}

\end{document}
